\subsection{2.5. Ataques Reales}\label{ataques-reales}

El estudio de incidentes reales constituye una fuente de conocimiento
imprescindible en el ámbito de la seguridad de contratos inteligentes. A
diferencia de los entornos de prueba, los ataques producidos en redes de
producción demuestran cómo las vulnerabilidades teóricas se traducen en
pérdidas económicas concretas e irreversibles.

En esta sección se analizan cuatro casos históricos representativos,
seleccionados por su relevancia técnica y su correspondencia directa con
cada una de las categorías de vulnerabilidades descritas en la sección
anterior (4.4. Vulnerabilidades de seguridad en contratos inteligentes
de Ethereum).

\subsubsection{2.5.1. Vulnerabilidad técnica de ejecución: The DAO
(2016)}\label{vulnerabilidad-tuxe9cnica-de-ejecuciuxf3n-the-dao-2016}

El ataque contra The DAO, producido en junio de 2016, representa el caso
más paradigmático de explotación de una vulnerabilidad de
\emph{reentrancy} en la historia de Ethereum {[}36{]}. The DAO era un
fondo de inversión descentralizado desplegado sobre Ethereum que había
recaudado aproximadamente 150 millones de dólares en ETH procedentes de
más de once mil participantes, convirtiéndose en uno de los mayores
proyectos de financiación colectiva hasta ese momento {[}37{]}.

La vulnerabilidad residía en la función withdraw() del contrato. El
flujo de ejecución enviaba los fondos en ETH al destinatario antes de
actualizar el balance interno correspondiente. Esta secuencia incorrecta
permitió a un atacante desplegar un contrato receptor cuya función de
\emph{fallback} re-invocaba withdraw() de forma recursiva, extrayendo
fondos de manera repetida mientras el saldo del contrato permanecía sin
modificar. El ataque se ejecutó en múltiples iteraciones dentro de una
misma cadena de llamadas, drenando aproximadamente 3,6 millones de ETH
{[}38{]}.

Las consecuencias del incidente trascendieron lo puramente económico. La
comunidad de Ethereum debatió durante semanas sobre cómo dar respuesta
al ataque, dado el carácter inmutable del protocolo. Finalmente, se
ejecutó un \emph{hard fork} en el bloque 1.920.000, el 20 de julio de
2016, que revirtió el estado de la cadena para recuperar los fondos. Una
fracción de la comunidad rechazó esta intervención por considerarla
contraria a los principios de inmutabilidad de la tecnología blockchain,
lo que dio lugar a la bifurcación conocida como Ethereum Classic
{[}37{]}.

Este caso puso de manifiesto que el patrón de diseño correcto para
gestionar transferencias es el denominado
\emph{checks-effects-interactions}: primero verificar condiciones,
después actualizar el estado del contrato y, por último, realizar
cualquier llamada externa. Su incumplimiento en The DAO, a pesar de la
aparente sencillez del principio, resultó en pérdidas valoradas en torno
a 60 millones de dólares al precio del momento {[}36{]}.

\subsubsection{2.5.2. Vulnerabilidad de control y privilegios: Poly
Network
(2021)}\label{vulnerabilidad-de-control-y-privilegios-poly-network-2021}

El ataque a Poly Network, ocurrido el 10 de agosto de 2021, ilustra con
especial claridad las consecuencias de un control de acceso deficiente
en contratos que gestionan operaciones críticas de alto valor {[}39{]}.
Poly Network es un protocolo de interoperabilidad \emph{cross-chain} que
conecta diversas redes blockchain, lo que amplía considerablemente su
superficie de ataque.

El vector de explotación se basó en una ausencia de restricciones en la
cadena de llamadas entre contratos del protocolo. La función
verifyHeaderAndExecuteTx(), presente en el contrato
EthCrossChainManager, permitía invocar internamente la función
PutCurEpochConPubKeyBytes() del contrato EthCrossChainData, encargada de
actualizar la clave pública del grupo de \emph{keepers} con permisos
para ejecutar transacciones \emph{cross-chain}. Al no existir ninguna
comprobación que restringiese qué direcciones podían realizar esta
invocación de forma indirecta, el atacante fue capaz de sustituir los
\emph{keepers} legítimos por una dirección bajo su propio control.

La operación se replicó de forma simultánea en tres redes: Ethereum, BNB
Chain y Polygon; lo que supuso un control inmediato sobre la práctica
totalidad de los fondos bloqueados en el protocolo, con un valor
estimado de 611 millones de dólares, convirtiéndose en el mayor robo de
la historia del ecosistema DeFi hasta esa fecha.

El desenlace del incidente fue inusual: el atacante devolvió la
totalidad de los fondos en los días posteriores, argumentando haber
actuado con el propósito de demostrar la vulnerabilidad. No obstante, el
caso evidencia que errores de diseño en el modelo de permisos
(independientemente de su recuperabilidad) pueden comprometer
completamente la integridad de un protocolo. Una correcta implementación
de control de acceso, mediante modificadores de autorización explícitos
y separación de privilegios entre contratos, habría impedido la escalada
de permisos aprovechada en este ataque {[}40{]}.

\subsubsection{2.5.3. Vulnerabilidad económica y dependencia del
entorno: bZx
(2020)}\label{vulnerabilidad-econuxf3mica-y-dependencia-del-entorno-bzx-2020}

Los ataques contra el protocolo bZx, producidos en febrero de 2020,
constituyen uno de los primeros ejemplos documentados de explotación
combinada de \emph{flash loans} y manipulación de oráculos de precio a
escala en el ecosistema DeFi. Su relevancia radica no solo en las
pérdidas económicas, sino en haber demostrado que la composabilidad de
los protocolos DeFi puede convertirse en un vector de ataque cuando los
componentes individuales no contemplan escenarios de manipulación
externa {[}41{]}.

En el primero de los dos incidentes, acaecido el 14 de febrero de 2020,
el atacante obtuvo un préstamo flash de 10.000 ETH a través del
protocolo dYdX sin necesidad de aportar colateral alguno. Con una parte
de estos fondos abrió simultáneamente una posición corta apalancada
sobre el par WBTC/ETH en bZx y utilizó el volumen restante para adquirir
WBTC en Uniswap, manipulando artificialmente el precio de mercado del
par en ese pool. El protocolo bZx, al depender de Uniswap como fuente de
precios en tiempo real, calculó incorrectamente la valoración de la
posición del atacante, quien la cerró obteniendo un beneficio aproximado
de 350.000 dólares una vez devuelto el préstamo flash, todo dentro de
una única transacción {[}42{]}.

El segundo ataque, cuatro días después, siguió una mecánica similar pero
orientada a la manipulación del precio del token sUSD en Kyber Network.
Al inflar artificialmente su cotización, el atacante logró depositar
colateral sobrevalorado en bZx y extraer en préstamo un volumen de
activos muy superior al que le habría correspondido según condiciones de
mercado normales {[}42{]}.

Ambos ataques pusieron de manifiesto que la seguridad de un protocolo
DeFi no puede analizarse de forma aislada: la dependencia de precios
obtenidos de fuentes únicas y manipulables en tiempo real constituye una
vulnerabilidad estructural. La adopción de oráculos descentralizados con
agregación de múltiples fuentes de precio (como los proporcionados por
protocolos del tipo Chainlink) y el uso de precios promediados en el
tiempo (TWAP, \emph{Time-Weighted Average Price}) son las principales
contramedidas frente a este patrón de ataque.

\subsubsection{2.5.4. Error lógico de negocio: Euler Finance
(2023)}\label{error-luxf3gico-de-negocio-euler-finance-2023}

El ataque a Euler Finance, ejecutado el 13 de marzo de 2023, representa
un ejemplo de elevada complejidad técnica dentro de la categoría de
errores lógicos, precisamente porque la vulnerabilidad no residía en un
patrón de codificación inseguro clásico, sino en una interacción no
prevista entre mecanismos financieros del propio protocolo {[}43{]}.

Euler Finance era un protocolo de préstamos descentralizado en Ethereum
que introducía el concepto de \emph{sub-cuentas} para gestionar
posiciones de colateral y deuda de forma independiente. La
vulnerabilidad se encontraba en la función donateToReserves(),
incorporada en una actualización anterior. Esta función permitía a un
usuario transferir sus propios activos a las reservas del protocolo,
pero no verificaba que dicha operación no dejase la cuenta del donante
en un estado de insolvencia. En condiciones normales, esta situación
habría sido detectada por el sistema de liquidación, pero el atacante
aprovechó el mecanismo de autoliquidación del protocolo de una forma no
contemplada en su diseño {[}44{]}.

El flujo del ataque comenzó con la obtención de un préstamo flash de 30
millones de DAI desde el protocolo Aave. Con estos fondos, el atacante
depositó colateral en Euler y, mediante operaciones de apalancamiento
repetidas, construyó posiciones de aproximadamente 195 millones de eDAI
(depósitos) y 200 millones de dDAI (deuda). A continuación, utilizó
donateToReserves() para crear intencionalmente una posición insolvente y
acto seguido activó la autoliquidación. El sistema de Euler, al calcular
el bonus de liquidación sobre una posición de ese tamaño, transfirió al
atacante los fondos de las reservas del protocolo en una cuantía muy
superior a la deuda original {[}45{]}.

La pérdida total ascendió a aproximadamente 197 millones de dólares. Sin
embargo, el atacante devolvió la práctica totalidad de los fondos
semanas después, tras una negociación directa con el equipo de Euler
mediante mensajes publicados en la cadena de bloques. El incidente
ilustra cómo la ausencia de validaciones sobre invariantes de negocio
(en este caso, que ninguna operación debería permitir dejar una cuenta
en estado de insolvencia) puede generar vulnerabilidades que escapan
tanto a la revisión manual del código como a las herramientas de
análisis estático-convencionales {[}44{]}.
